\documentclass[10pt,a4paper]{amsart}
\usepackage[margin=19mm]{geometry}
\usepackage{amsmath,amssymb,mathtools}
\usepackage[scr=rsfs,cal=euler]{mathalfa}
\usepackage{xfrac}
\usepackage[unicode,hidelinks,bookmarks=false]{hyperref}
\newcommand{\ie}{i.e.\ }
\newcommand{\cf}{cf.\ }
\newcommand{\resp}{respectively\ }
\newcommand{\Subsection}{\subsection}
\providecommand{\nobreakdash}{\nobreak\hbox{-}}
\setlength{\emergencystretch}{2em}
\multlinegap0pt
\renewcommand{\labelenumi}{(\arabic{enumi})}
\usepackage{cleveref}
\newtheorem{thm}{Theorem}[section]
\newtheorem{lm}[thm]{Lemma}
\newtheorem{res}[thm]{Result}
\newtheorem{crl}[thm]{Corollary}
\newtheorem{prop}[thm]{Proposition}
\newtheorem{const}[thm]{Construction}
\newtheorem{obs}[thm]{Observation}
\newtheorem{rmk}[thm]{Remark}
\newtheorem{df}[thm]{Definition}
\newtheorem{nota}[thm]{Notation}
\newtheorem{ex}[thm]{Example}
\newcommand{\eps}{\varepsilon}
\renewcommand{\phi}{\varphi}
\newcommand{\RR}{\mathbb R}
\newcommand{\FF}{\mathbb F}
	\newcommand{\fq}{\FF_q}
	\newcommand{\fp}{\FF_p}
  \newcommand{\fqi}{\fq \cup \{\infty\}}
\newcommand{\CC}{\mathbb C}
\newcommand{\QQ}{\mathbb Q}
\newcommand{\ZZ}{\mathbb Z}
\newcommand{\NN}{\mathbb N}
\renewcommand{\leq}{\leqslant}
\renewcommand{\geq}{\geqslant}
\newcommand{\vspan}[1]{\left \langle #1 \right \rangle}
\newcommand{\vspann}[2]{\left \langle #1 \, \| \, #2 \right \rangle}
\newcommand{\set}[1]{ \left \{ #1 \right \} }
\newcommand{\sett}[2]{ \left\{ #1 \, \, || \, \, #2 \right \} }
\newcommand{\gauss}[2]{\genfrac{[}{]}{0pt}{}{#1}{#2}_q}
\newcommand{\one}{\mathbf 1}
\newcommand{\zero}{\mathbf 0}
\newcommand{\floor}[1]{\left \lfloor #1 \right \rfloor}
\newcommand{\ceil}[1]{\left \lceil #1 \right \rceil}
\newcommand{\ma}{\mathcal A}
\newcommand{\mb}{\mathcal B}
\newcommand{\mc}{\mathcal C}
\newcommand{\md}{\mathcal D}
\newcommand{\me}{\mathcal E}
\newcommand{\mf}{\mathcal F}
\newcommand{\mg}{\mathcal G}
\newcommand{\mh}{\mathcal H}
\newcommand{\mi}{\textnormal I}
\newcommand{\mj}{\mathcal J}
\newcommand{\mk}{\mathcal K}
\newcommand{\ml}{\mathcal L}
\newcommand{\mm}{\mathcal M}
\newcommand{\mn}{\mathcal N}
\newcommand{\mo}{\mathcal O}
\renewcommand{\mp}{\mathcal P}
\newcommand{\mq}{\mathcal Q}
\newcommand{\mr}{\mathcal R}
\newcommand{\ms}{\mathcal S}
\newcommand{\mt}{\mathcal T}
\newcommand{\mv}{\mathcal V}
\newcommand{\mw}{\mathcal W}
\newcommand{\mx}{\mathcal X}
\newcommand{\my}{\mathcal Y}
\newcommand{\mz}{\mathcal Z}
\newcommand{\sam}[1]{\textcolor{purple}{#1}}
 \DeclareMathOperator{\supp}{supp}
 \DeclareMathOperator{\wt}{wt}
 \DeclareMathOperator{\rk}{rk}
 \DeclareMathOperator{\Tr}{Tr}
    \newcommand{\tr}{\Tr}
 \DeclareMathOperator{\PG}{PG}
    \newcommand{\pg}{\PG}
 \DeclareMathOperator{\AG}{AG}
    \newcommand{\ag}{\AG}
 \DeclareMathOperator{\PGL}{PGL}
    \newcommand{\pgl}{\PGL}
 \DeclareMathOperator{\PSL}{PSL}
    \newcommand{\psl}{\PSL}
 \DeclareMathOperator{\PGamL}{P\Gamma L}
 \DeclareMathOperator{\GL}{GL}
 \DeclareMathOperator{\GamL}{\Gamma L}
 \DeclareMathOperator{\PGO}{PGO}
  \newcommand{\pgo}{\PGO}
 \DeclareMathOperator{\id}{id}
 \DeclareMathOperator{\proj}{proj}
 \DeclareMathOperator{\NO}{NO}
  \newcommand{\nono}{\overline{\NO^+}(8,2)}
 \DeclareMathOperator{\Sym}{Sym}
 \newcommand{\BP}{\boldsymbol P}
 \DeclareMathOperator{\SRG}{SRG}
 \DeclareMathOperator{\Aut}{Aut}
 \DeclareMathOperator{\Johnson}{J}
 \DeclareMathOperator{\GO}{GO}
 \DeclareMathOperator{\OS}{OS}
 \DeclareMathOperator{\Alt}{Alt}
 \DeclareMathOperator{\ASL}{ASL}
 \DeclareMathOperator{\Inv}{Inv}
\begin{document}
\title[A sporadic strongly regular graph with parameters (120,56,28]{A sporadic strongly regular graph with parameters $(120,56,28,24)$ from a primitive action of the symmetric group on $7$ elements}
\author{Jan De Beule [1]; Sam Adriaensen, Robert F. Bailey, Jan De Beule, Morgan Rodgers}
\date{}
\maketitle
\noindent\emph{Source and licence.} A sporadic strongly regular graph with parameters $(120,56,28,24)$ from a primitive action of the symmetric group on $7$ elements. Version arXiv:2606.10652v1 (CC BY 4.0). This material is distributed under the Creative Commons Attribution 4.0 International licence, http://creativecommons.org/licenses/by/4.0/, as recorded in the arXiv OAI licence metadata for this exact version and shown on the abstract page https://arxiv.org/abs/2606.10652v1. Full rights notice: licenses/arxiv-2606.10652-CC-BY-4.0.txt.\par\smallskip
\noindent\textit{Editorial scope.} Counted unit: the original \texttt{lm} environment \texttt{-} at source lines 531--538 together with its complete proof at lines 540--551; the delivered document keeps source lines 360--552 so that every referenced label and every citation resolves inside it. Native editable source is the paper's own concatenated TeX; the AMS wrapper is editorial formatting only. No publisher class, upstream script or unrelated artwork is redistributed.
\begin{align}
 \label{Eq:CP}
 V = \sett{(x_1, \ldots, x_9) \in \FF_2^9}{x_1 + \ldots + x_9 = 0} \cong \FF_2^8,
 && f(X) = \sum_{1 \leq i < j \leq 9} X_i X_j.
\end{align}
The \emph{weight} of a vector is defined as the number of coordinate positions with a non-zero entry.
If $x \in V$ has weight $n$, then over $\FF_2$,
\[
 f(x) = \binom n2 = (n-1) \frac n2 = \frac n2.
\]
Thus, $\vspan x$ is an isotropic point if and only if $x$ has weight 4 or 8.
Note that the points of weight 8 form an ovoid.

Let $A = \{1, \dots, 9\}$.
We can identify each vector $x \in \FF_2^9$ with the set $\sett{a \in A}{x_a = 1}$.
Cameron and Praeger \cite{Cameron:Praeger} observed that a set $\mb \subset \binom A4$ constitutes a $3$-$(8,4,1)$ design on $A \setminus \{a\}$ if and only if $\mb \cup \{ A \setminus \{a\}\}$ are the points of a generator in $\mq^+(7,2)$.
This in particular yields a one-to-one correspondence between spreads of $\mq^+(7,2)$ and $\OS(4,8)$'s.

The group $\Sym_9$ naturally acts on $V$ by coordinate permutation.
Coordinate permutation does not alter the quadratic form $f$, so we can see $\Sym_9$ as a subgroup of $\PGO^+(8,2)$.
Its action on the spreads of $\mq^+(7,2)$ coincides with the action of $\Sym_9$ on the $\OS(4,8)$'s induced by the natural action of $\Sym_9$ on the 9-element ground set.
Since every system of generators has the same number of spreads, and the action of $\Sym_9$ has orbits of sizes $240$ and $1680$ on the spreads of $\mq^+(7,2)$, it cannot stabilise the systems of generators.
Thus, its subgroup stabilising the systems of generators must be of index 2, hence be $\Alt_9$.
We know the orbits of $\Alt_9$ on the $\OS(4,8)$'s.
Hence, in each system of generators, the orbits of $\Alt_9$ on the spreads have sizes $120$, corresponding to $\OS(4,8)$'s of type (1), and $840$, corresponding to $\OS(4,8)$'s of type (2).
Moreover, as we argued, the subgroup of $\Alt_9$ that setwise stabilises a pair of coordinate positions has the same orbits on the spreads of $\mq^+(7,2)$.

\subsection{The association scheme}
 \label{Sec:Assoc:Spreads}

We will now present an association scheme defined on $120$ $\OS(4,8)$'s of type (1) that form an orbit under the action of $\Alt_9$.
However, we will rephrase everything to the language of spreads of $\mq^+(7,2)$.
We start with the following simple observation.

\begin{lm}
 Let $\ms$ be a spread of $\mq^+(7,2)$, and $\Sigma$ a generator from the same system as the generators of $\ms$.
 Then either $\Sigma \in \ms$, or $\Sigma$ intersects five generators of $\ms$ in a line, and is disjoint from the other four.
\end{lm}

\begin{proof}
 Each generator $\Sigma_i$ of $\ms$ lives in the same system of generators as $\Sigma$, hence intersects it in a subspace of dimension $-1$, $1$, or $3$.
 Moreover, every point of $\Sigma$ is on a unique generator of $\ms$.
 Hence, if $\Sigma \notin \ms$, the non-empty intersections of $\Sigma$ with the elements of $\ms$ partition the points of $\Sigma$ into lines.
 Since $\Sigma$ has 15 points, and a line has three points, such a partition contains five lines.
\end{proof}

Let $V$ and $f$ be as in (\ref{Eq:CP}).
Let $Q_i$ denote the point whose coordinate vector has a single 0 in position $i$.
Consider $\Alt_9$ acting on $V$ by coordinate permutation, and let $G^*$ denote the subgroup of $\Alt_9$ that setwise stabilises coordinates 1 and 2.
Then we saw that in each system of generators of $\mq^+(7,2)$, $G^*$ has two orbits of sizes $120$ and $840$ on the spreads.
Let $\mx$ be such an orbit of size $120$.

\begin{prop}[GAP]
 \label{Prop:Gen Trans}
 The action of $G^*$ on $\mx$ is generously transitive and has rank $7$.
\end{prop}

Next, we introduce two quantities that are invariant under the action of $G^*$.
Given a spread $\ms \in \mx$, let $\ms_i$ denote the generator in $\ms$ containing the point $Q_i$.
Given two spreads, $\ms$ and $\mt$, the number
\[
 \Inv_1(\ms,\mt) = \dim(\ms_1 \cap \mt_2) + \dim(\ms_2 \cap \mt_1)
\]
is easily checked to be $G^*$-invariant.
We define $\Inv_2(\ms,\mt)$ to be the number of generators in $\{\mt_3, \dots, \mt_9\}$ that intersect both $\ms_1$ and $\ms_2$ in a line.
This is also easily checked to be $G^*$-invariant.
Note that by $G^*$ acting generously transitively, we know that $\Inv_i(\ms,\mt) = \Inv_i(\mt,\ms)$ for $i = 1,2$.

We define relations $R_0, \dots, R_6 \subset \mx^2$, where $R_0$ is simply the identity relation, and relations $R_1, \dots, R_6$ are defined using the Table \ref{tab:Assoc1}.

\begin{table}[h!]
 \centering
 \begin{tabular}{|l||c | c | c | c|}
  \hline
  & $\Inv_2(\ms,\mt) = 0$ & $\Inv_2(\ms,\mt) = 1$ & $\Inv_2(\ms,\mt) = 2$ & $\Inv_2(\ms,\mt) = 3$  \\ \hline \hline
  $\Inv_1(\ms,\mt) = -2$ & $R_2$ & & $R_3$ & \\ \hline
  $\Inv_1(\ms,\mt) = 0$ & & $R_5$ & & $R_6$ \\ \hline
  $\Inv_1(\ms,\mt) = 2$ & $R_1$ & & $R_4$ & \\ \hline
 \end{tabular}
 \caption{Relations of the association scheme on $120$ spreads of $\mq^+(7,2)$.}
 \label{tab:Assoc1}
\end{table}

\begin{thm}[GAP]
 \label{Thm:Assoc:Spreads}
 The 6-class symmetric association scheme on $\mx$, formed by the orbitals of the generously transitive action of $G^*$, is given by the relations $R_0, \dots, R_6$ as described in Table \ref{tab:Assoc1}.
 The matrix of eigenvalues is given by
 \[
  \BP = \begin{pmatrix}
1 & 7 & 14 & 21 & 21 & 42 & 14\\
1 & 2 & 4 & -9 & 6 & -3 & -1\\
1 & -1 & -2 & 3 & 3 & -6 & 2\\
1 & -4 & -2 & -3 & 0 & 9 & -1\\
1 & 3 & -2 & 3 & -1 & 2 & -6\\
1 & -2 & 8 & 3 & -6 & -3 & -1\\
1 & 3 & -2 & -3 & -7 & 2 & 6\\
\end{pmatrix}
 \]
 There are two possible fusions of the relations that yield an $\SRG(120,56,28,24)$:
\begin{align*}
 \BP (0,0,0,0,0,1,1)^\top
 = (56,-4,8,-4,-4,-4,8)^\top, &&
 \BP (0,0,0,1,1,0,1)^\top
 = (56,-4,-4,-4,8,-4,-4)^\top.
\end{align*}
 The graph given by $R_5 \cup R_6$ is isomorphic to $\nono$.
 The graph $\Gamma^*$ given by $R_3 \cup R_4 \cup R_6$ has $G^*$ as its full automorphism group.
\end{thm}



\begin{rmk}[GAP]
 \label{Rmk:extra invariant}
Another natural $G^*$-invariant is of course checking how many generators $\ms$ and $\mt$ have in common, and we can even split this into 2 $G^*$-invariants: $|\{\ms_1, \ms_2\} \cap \{\mt_1, \mt_2\}|$ and $|\{\ms_3, \dots, \ms_9\} \cap \{\mt_3, \dots, \mt_9\}|$.
If we replace $\Inv_1$ or $\Inv_2$ by this invariant, we can no longer distinguish all the relations of the association scheme.
However, this invariant gives useful information.
We describe it in Table \ref{tab:extra invariant}.
 Note that two spreads $\ms$ and $\mt$ in $\mx$ share at most one generator.
 Moreover, we can define $\nono$ as the graph on $\mx$ where two spreads are adjacent if they share no generator.

\begin{table}[h!]
 \centering
 \begin{tabular}{|c||c|c|}
 \hline
 & $|\{\ms_3, \dots, \ms_9\} \cap \{\mt_3, \dots, \mt_9\}| = 0$ & $|\{\ms_3, \dots, \ms_9\} \cap \{\mt_3, \dots, \mt_9\}| = 1$ \\ \hline \hline
 $|\{\ms_1, \ms_2\} \cap \{\mt_1, \mt_2\}| = 0$ & $R_5$, $R_6$ & $R_1$, $R_3$, $R_4$ \\ \hline
 $|\{\ms_1, \ms_2\} \cap \{\mt_1, \mt_2\}| = 1$ & $R_2$ & \\ \hline
 \end{tabular}
 \caption{An extra invariant on the association scheme.}
 \label{tab:extra invariant}
\end{table}
\end{rmk}

\section{Ovoids and anisotropic points \texorpdfstring{of $\mq^+(7,2)$}{}}

If we represent $\mq^+(7,2)$ using the vector space and quadratic form from (\ref{Eq:CP}), then the points $Q_1, \dots, Q_9$ whose coordinate vector has weight 8 form an ovoid $\mo$ of $\mq^+(7,2)$.
Clearly, the group $\Sym_9 \leq \PGO^+(8,2) \leq \PGL(8,2)$ of coordinate permutations stabilises $\mo$ setwise.
The ovoid is a frame and since $\PGL(8,2)$ acts sharply transitively on ordered frames, we see that $\Sym_9$ is the full setwise stabiliser of $\mo$ in $\PGO^+(8,2)$.
Moreover, $\Alt_9$ is the setwise stabiliser of $\mo$ that also stabilises the systems of generators.

All ovoids of $\mq^+(7,2)$ are isomorphic, see e.g.\ \cite{Cameron:Praeger}.
Thus, we can give a coordinate-free definition of the group $G^*$ from the previous section, by letting it be the setwise stabiliser of any ovoid of $\mq^+(7,2)$ that stabilises the systems of generators.
Using the principle of triality, we can translate the results of the previous section to the following proposition.

\begin{prop}
 \label{Prop:by triality}
 Let $\ms$ be a spread of $\mq^+(7,2)$, and let $\Sigma_1, \Sigma_2 \in \ms$.
 Let $G^*$ be the setwise stabiliser of $\{\Sigma_1, \Sigma_2\}$ and $\ms \setminus \{\Sigma_1, \Sigma_2\}$ in $\PGO^+(8,2)$.
 Then $G^*$ has two orbits of sizes $120$ and $840$ on the ovoids of $\mq^+(7,2)$.
 Let $\my$ denote the orbit of size $120$.
 The action of $G^*$ on $\my$ is generously transitive of rank $7$.
 Moreover, distinct ovoids in $\my$ share at most one point.
\end{prop}

This means that any ovoid $\mo \in \my$ is uniquely determined by the points $\mo \cap \Sigma_1$ and $\mo \cap \Sigma_2$.
On the other hand, there are 15 choices for a point $P_1 \in \Sigma_1$, and 8 choices for a point $P_2 \in \Sigma_2 \setminus P_1^\perp$, which yields a total of $15 \cdot 8 = 120 = |\my|$ pairs of non-collinear points.
This gives a one-to-one correspondence.
On the other hand, since the vector space underlying $\pg(7,2)$ is a direct sum of the subspaces corresponding to $\Sigma_1$ and $\Sigma_2$, given a point $P = \vspan x \notin \Sigma_1 \cup \Sigma_2$, there is a unique way to write $x$ as a linear combination of a vector in $\Sigma_1$ and $\Sigma_2$ seen as vector subspaces.
This means that there is a unique line $\ell_P$ through $P$ that intersects both $\Sigma_1$ and $\Sigma_2$.
Note that given two points $P_1 \in \Sigma_1$ and $P_2 \in \Sigma_2$, the line $P_1 P_2$ contains one extra point $P_3$, and this point belongs to $\mq^+(7,2)$ if and only if $P_1 \perp P_2$.
Therefore, there is a one-to-one correspondence between anisotropic points and pairs of non-collinear points in $\Sigma_1, \Sigma_2$.

To summarise, for each anisotropic point $P$, there is a unique line $\ell_P$ through $P$ that intersects $\Sigma_1$ and $\Sigma_2$, and a unique ovoid $\mo_P \in \my$ that contains $\Sigma_1 \cap \ell_P$ and $\Sigma_2 \cap \ell_P$.
This gives a one-to-one correspondence between anisotropic points and the ovoids of $\my$, and this correspondence is $G^*$-invariant.

In particular, this allows us to translate the association scheme from Section \ref{Sec:Assoc:Spreads} to an association scheme on the anisotropic points.
The relations of this scheme are best described by not only considering the anisotropic points, but also their corresponding ovoids.
We will denote the ovoid of $\my$ corresponding to the anisotropic point $P$ by $\mo_P$.
Given anisotropic points $P$ and $Q$, a useful $G^*$-invariant property is the number of points in $\mo_P$ orthogonal to $Q$.
There are only two options.


\begin{lm}
 Let $\mo$ be an ovoid of $\mq^+(7,2)$, and $P$ an anisotropic point.
 Then one of the following holds:
 \begin{enumerate}
  \item Exactly one line through $P$ intersects $\mo$ in 2 points, and $|P^\perp \cap \mo| = 7$.
  \item No line through $P$ intersects $\mo$ in 2 points, and $|P^\perp \cap \mo| = 3$.
 \end{enumerate}
\end{lm}

\begin{proof}
 We can use the representation of (\ref{Eq:CP}) for $V$ and $f$.
 Let $e_1, \dots, e_9$ denote the standard basis vectors of $\FF_2^9$.
 We can apply a coordinate transformation such that $\mo$ consists of the weight 8 points.
 Let $P$ be an anisotropic point with coordinate vector $v$.
 If $v$ has weight $2$, then $v = e_i + e_j$.
 Given a vector $w = \one + e_k$ of weight $8$, $v+w$ has weight $8$ if $k \in \{i,j\}$, and weight $6$ otherwise.
 Hence, $P$ is orthogonal the points $\vspan{\one + e_k}$ with $k \notin \{i,j\}$, and $P$, $\vspan{\one + e_i}$, $\vspan{\one + e_j}$ lie on the same line.
 If $v$ has weight $6$, then $v = \one + e_i + e_j + e_k$.
 Given a vector $w = \one + e_h$ of weight 8, $v+w$ has weight 2 if $h \in \{i,j,k\}$, and weight 4 otherwise.
 Thus, no two points of weight 8 lie on the same line through $P$, and there are exactly 3 points of $\mo$ which lie on a tangent line to $\mq^+(7,2)$ through $P$, hence are orthogonal to $P$.
\end{proof}


\begin{thebibliography}{99}
\bibitem[CP91]{Cameron:Praeger} P.~J. Cameron and C.~E. Praeger. \newblock Partitioning into {S}teiner systems. \newblock In {\em Combinatorics '88, {V}ol.\ 1 ({R}avello, 1988)}, Res. Lecture Notes Math., pages 61--71. Mediterranean, Rende, 1991.
\end{thebibliography}
\end{document}
